\chapter{Introduction}

\section{Background}
Reinforced concrete (RC) frame buildings are among the most common building systems in urban and semi-urban areas of Nepal. Rapid urbanization, limited plots, and the demand for residential floor area have encouraged widespread construction of low-rise RC buildings, often with substantial involvement of owner-builders, local technicians, and masons. Simplified Nepalese provisions have therefore played an important role in providing minimum member sizes and reinforcement arrangements where a complete project-specific structural design may not always be available.

The evolution of these simplified provisions must be distinguished clearly. NBC 205:1994 introduced the Mandatory Rules of Thumb (MRT), while the Draft Final NBC 205:2012 and NBC 205:2024 documents provide Ready-to-Use Detailing (RUD) guidance for low-rise RC buildings without masonry infill \cite{nbc205_1994,nbc205_2012,nbc205_2024}. In this thesis, the representative 2012 building refers to the building designed using the Draft Final NBC 205:2012 ready-to-use provisions. NBC 105:2025 provides the seismic demand framework, including seismic zone factors, soil classes, equivalent-static provisions, and response-spectrum definitions \cite{nbc105_2025}.

Nepal lies in a highly seismic region because of the continuing convergence of the Indian and Eurasian tectonic plates. The 1934 Nepal--Bihar earthquake, the 2015 Gorkha earthquake, and other damaging events have demonstrated the vulnerability of Nepalese building stock. Previous research has shown that construction quality, ductile detailing, code compliance, and structural configuration all influence earthquake performance \cite{thapa2013psha,parajuli2015gorkha,chaulagain2013seismic}.

Building irregularity is particularly important because it can disturb the distribution of mass, stiffness, strength, and lateral-force resistance. L-, T-, U-shaped, and non-parallel configurations may produce re-entrant corners, eccentricity between the centre of mass and centre of rigidity, CM--CR eccentricity, concentrated drift, and non-uniform hinge formation. Height irregularity can also change the vertical distribution of stiffness, mass, and seismic demand. Consequently, a design that satisfies prescriptive member-size and reinforcement rules may still require detailed analysis when the building departs from the regular configurations assumed by simplified guidance \cite{cimellaro2014bidirectional,destefano2014pushover,jadhav2025study}.

Performance-based seismic evaluation provides a rational way to assess these effects beyond basic strength checks. Equivalent static and response spectrum analyses describe elastic force and deformation demand, whereas nonlinear static pushover analysis provides capacity curves, performance points, and hinge-state information. These indicators should be interpreted together because peak strength, demand-state base shear, displacement demand, torsion, and damage progression represent different aspects of structural performance \cite{atc1996atc40,fema2005fema440,asce2017asce41}.

The equivalent-static method distributes an idealized lateral force over the building height. Response-spectrum analysis combines modal demands derived from the applicable seismic spectrum. Nonlinear static pushover analysis traces an idealized capacity response under increasing lateral load. These analytical concepts are developed in the methodology; the numerical findings are presented in Chapter 4.

For pushover interpretation, the embedded performance-point response, the first global peak, the maximum recorded displacement, and the available final-step hinge evidence are treated as distinct response quantities rather than interchangeable indicators. Final-step hinge evidence describes only the recorded final analysis step and is not used to infer a hinge state at the performance point. A performance-point hinge state is reported only when it is directly present in a solved hinge output and is not interpolated from final-step evidence.

The research follows four linked phases. Phase I reviews a representative building detailed using Draft Final NBC 205:2012 using the available concrete design-check results. Phase II establishes the NBC 205:2024 Regular building as the performance-based reference model. Phase III evaluates six plan-irregular 2024 buildings and a separately defined four-storey-plus-stair-cover height-extension model. Phase IV synthesizes the force, deformation, CM--CR separation, capacity, performance-point, hinge, and failure evidence to identify practical applicability limits of NBC 205 RUD provisions.

The study uses ETABS analysis results for structural response assessment. The principal response parameters are storey shear, inter-storey drift, CM--CR separation, peak pushover capacity, performance-point response, final-step hinge observations, and concrete design-check observations. The 2012 model is retained as a historical design-check reference, while the nonlinear pushover and performance-point comparisons are based on the 2024 analytical models.

This chapter presents the background, problem statement, research questions, objectives, and significance of the study. The subsequent chapters review relevant literature, describe the analytical methodology, discuss the analysis results, and present conclusions and recommendations.

Nepal is located in a seismically active region created by the continuing convergence of the Indian plate beneath the Eurasian plate. Destructive earthquakes, including the 1934 Nepal--Bihar event and the 2015 Gorkha earthquake, have repeatedly demonstrated the vulnerability of the national building stock. Low-rise residential RC buildings require particular attention because they are frequently constructed under variable levels of engineering input, workmanship, and quality control \cite{thapa2013psha,parajuli2015gorkha}.

Simplified design and detailing documents were developed to improve minimum safety in this construction environment. NBC 205:1994 provided Mandatory Rules of Thumb, while the Draft Final NBC 205:2012 and NBC 205:2024 documents provide ready-to-use detailing guidance for low-rise RC buildings without masonry infill \cite{nbc205_1994,nbc205_2012,nbc205_2024}. The 2012 document represents the basis of many buildings constructed before and around the 2015 Gorkha earthquake, whereas the 2024 edition provides the revised detailing basis evaluated in the present study.

Among the applicability restrictions of NBC 205:2024, buildings shall not have any irregularity, and the maximum structural height is 12 m or three storeys, whichever is less, measured from the level of lateral restraint. Within the 12 m height limit, a smaller additional storey may be provided only when its plan area does not exceed 25\% of a typical floor and its storey height is limited to 2.4 m; if this allowance is exceeded, it is treated as an additional storey. These restrictions define the prescriptive RUD eligibility boundary considered in this study and are distinct from the analytical seismic-response checks performed under NBC 105:2025 \cite{nbc205_2024}.

The seismic design basis is provided separately by NBC 105:2025, which defines seismic zone factors, soil classifications, response spectra, and related analysis requirements \cite{nbc105_2025}. The distinction is important: NBC 205 provides simplified member sizing and reinforcement detailing, whereas NBC 105:2025 defines the seismic demand applied in the analytical models.


\begin{figure}[!htbp]
\centering
\begin{tikzpicture}[>=Latex,node distance=1.2cm]
\node[draw,rounded corners,align=center,text width=0.24\textwidth,minimum height=1.15cm] (a) {\textbf{NBC 205:1994}\\Mandatory Rules of Thumb (MRT)};
\node[draw,rounded corners,align=center,text width=0.28\textwidth,minimum height=1.15cm,right=of a] (b) {\textbf{Draft Final NBC 205:2012}\\Ready-to-Use detailing reference};
\node[draw,rounded corners,align=center,text width=0.25\textwidth,minimum height=1.15cm,right=of b] (c) {\textbf{NBC 205:2024}\\Ready-to-Use Detailing (RUD)};
\draw[->,thick] (a) -- (b);
\draw[->,thick] (b) -- (c);
\end{tikzpicture}
\caption{Evolution of NBC 205 provisions.}
\label{fig:nbc205_evolution}
\end{figure}

The timeline presented in Figure~\ref{fig:nbc205_evolution} shows the progression from the earlier Mandatory Rules of Thumb to the 2012 ready-to-use reference and the current NBC 205:2024 RUD provisions. The sequence is used to distinguish the historical design basis from the current detailing framework assessed in this thesis.


\begin{figure}[!htbp]
\centering
\includegraphics[width=0.90\textwidth]{figures/ch1/nbc105_2025_seismic_zone_map.png}
\caption{Seismic zoning of Nepal (Source: NBC 105:2025 \cite{nbc105_2025}).}
\label{fig:nepal_seismic_zoning}
\end{figure}

The seismic-zoning map in Figure~\ref{fig:nepal_seismic_zoning} places the study within the national seismic-demand framework. The zone factors used in the analytical cases represent the code-defined regional demand levels and are applied together with the selected soil classes in the subsequent response evaluation \cite{nbc105_2025}.

\FloatBarrier


For the inter-storey drift check used in this study, NBC 105:2025 requires the ratio of inter-storey deflection to the corresponding storey height not to exceed 0.025 at the ultimate limit state and 0.006 at the serviceability limit state. These limits are used as analytical response checks and do not by themselves establish eligibility under the prescriptive scope of NBC 205:2024 \cite{nbc105_2025}.

NBC 105:2025 defines torsional irregularity using the variation in horizontal displacement across the ends of a storey. Torsional irregularity exists when the maximum horizontal displacement at one end of the storey is more than 1.5 times the minimum displacement at the far end in the same direction, while a ratio greater than 2.5 constitutes extreme torsional irregularity \cite{nbc105_2025}. In the comparative results, CM--CR separation is treated as an indicator of mass--stiffness eccentricity, rather than as the storey-end displacement ratio used for formal torsional-irregularity classification.

Previous studies on Nepalese RC buildings have shown that inadequate detailing, poor construction quality, weak-column strong-beam behaviour, and irregular drift concentration can increase vulnerability \cite{chaulagain2013seismic}. Research on buildings outside MRT geometric limitations has also shown that slenderness and plan aspect ratio affect seismic performance \cite{poudel2017seismic}. More recent work on NBC 205:2024 RUD buildings indicates that acceptable performance is possible when material quality, detailing, and construction supervision are maintained \cite{poudel2025seismic}.

The research flow separates the problem into four linked stages. First, a representative NBC 205:2012 building is reviewed using the available concrete design-check results. Second, a regular NBC 205:2024 building within the guideline's three-storey applicability range is evaluated as the performance-based benchmark. Third, selected plan-irregular configurations are analysed for comparison with that benchmark, while a separate four-storey-plus-stair-cover model represents a height-extension case. Fourth, the combined results are interpreted to identify practical applicability limits and cases that require project-specific structural assessment.

All regular and plan-irregular models used in the primary seven-model comparison have three full storeys plus a stair cover. Only the separate four-storey height-extension model has four full storeys plus a stair cover. This distinction is maintained consistently throughout the methodology and results.

The study therefore examines whether the NBC 205:2024 Regular benchmark provides an acceptable reference response and whether irregular configurations remain compatible with the assumptions of ready-to-use detailing. The 2012 model is evaluated through its concrete design-check results, while the pushover-capacity comparison concerns the 2024 models.

\section{Statement of the Problem}
Low-rise RC buildings in Nepal are often designed or constructed using simplified NBC 205 provisions because these provisions are accessible and practical for ordinary residential construction. However, compliance with prescriptive member dimensions and reinforcement details does not by itself demonstrate acceptable seismic performance under the demand defined by NBC 105:2025.

A representative existing building designed using Draft Final NBC 205:2012 provides a historical reference for assessing how earlier ready-to-use practice responds to the current NBC 105:2025 seismic-demand framework. Its evaluation is based on the available concrete design-check results. A direct nonlinear pushover comparison of the 2012 and 2024 models is outside the completed scope of the present study and is not used in the current conclusions.

The principal analytical problem is whether the NBC 205:2024 Regular model provides an acceptable low-rise benchmark and whether the same ready-to-use detailing concept can provide satisfactory seismic performance when applied to selected plan-irregular configurations. Re-entrant corners, asymmetric plans, non-parallel framing and additional height can alter load paths, eccentricity, torsion, displacement demand and nonlinear response in ways that simplified detailing rules alone do not describe.

Accordingly, this study evaluates the regular and selected irregular low-rise RC buildings under NBC 105:2025 and examines the practical applicability of NBC 205:2024 RUD provisions to the selected irregular configurations.

Analytical response and prescriptive-scope eligibility are treated as separate questions. A model may satisfy an available NBC 105:2025 response check while remaining outside the geometric or use restrictions of NBC 205:2024; conversely, being within the RUD scope does not remove the need for correct materials, detailing, construction, and seismic-demand evaluation. The study is therefore an applicability-boundary assessment rather than a code-certification exercise \cite{nbc205_2024,nbc105_2025}.

\section{Research Questions}
\begin{enumerate}[leftmargin=1.0cm,itemsep=0.30em,topsep=0.25em,parsep=0pt]
\item How does the representative low-rise RC building designed according to Draft Final NBC 205:2012 perform under the seismic demand specified by NBC 105:2025?
\item How does the regular low-rise RC building designed according to NBC 205:2024 RUD provisions perform under the seismic demand specified by NBC 105:2025?
\item How do the selected irregular low-rise RC buildings perform when the same NBC 205:2024 RUD provisions used for the Regular building are applied to them?
\item How do plan configuration and CM--CR separation affect the selected response measures: storey shear, inter-storey drift, pushover response, performance-point response, and hinge evidence?
\item What does the observed analytical response imply for the practical applicability boundary of NBC 205:2024 RUD provisions for the selected irregular configurations?
\item Is the Geometric Irregularity Index (GII) sufficiently representative of the observed differences in seismic response among the investigated plan configurations?
\end{enumerate}

\section{Objectives of the Study}
\noindent\textbf{General Objective:} To evaluate the seismic performance of low-rise reinforced concrete buildings designed using NBC 205:2024 Ready-to-Use Detailing (RUD) provisions and investigate the response and practical applicability boundary of selected irregular building configurations under NBC 105:2025 seismic demand.

\noindent\textbf{Specific Objectives:}
\begin{enumerate}[leftmargin=1.0cm,itemsep=0.25em,topsep=0.25em,parsep=0pt]
\item To evaluate the seismic response of a representative low-rise RC building designed according to Draft Final NBC 205:2012 using the seismic demand specified by NBC 105:2025 and examine the available concrete design-check results.
\item To evaluate the seismic response of a regular low-rise RC building designed according to NBC 205:2024 RUD provisions and establish it as the reference model.
\item To evaluate and compare the seismic response of selected irregular low-rise RC buildings designed using the same NBC 205:2024 RUD provisions.
\item To investigate the effects of plan configuration and CM--CR separation on storey shear, inter-storey drift, pushover response, performance-point response, and the available hinge evidence.
\item To examine the practical applicability boundary of NBC 205:2024 RUD provisions for the selected irregular configurations based on the observed analytical response.
\item To examine the relationship between GII and selected seismic response parameters.
\end{enumerate}

\section{Significance of the Study}
The provisions provided by NBC 205:2024 Ready-to-Use Detailing (RUD) suggest a simple approach for designing and detailing low-rise reinforced concrete buildings. Such provisions are useful because they can make the design process easier for buildings that satisfy the conditions specified by the guideline. However, the RUD provisions are intended for regular building configurations, while buildings constructed in practice may have different types of plan irregularity.

This study examines how selected plan-irregular buildings respond when they are detailed using the same NBC 205:2024 RUD provisions as the Regular reference building. No separate or specially modified detailing rules are introduced for the irregular configurations; their seismic response is evaluated under the same NBC 105:2025 demand framework used for the Regular model.

The study considers Regular, L-shaped, L1-shaped, L2-shaped, T-shaped, U-shaped, and Non-Parallel building configurations. Their responses are compared using storey shear, inter-storey drift, CM--CR separation, pushover capacity, performance points, and final-step hinge observations. Considering several response parameters provides a better basis for interpreting comparative behaviour than relying on only one parameter.

The comparison also provides a direct basis for judging whether any of the investigated irregular configurations remain close to the Regular reference across the selected response measures. Such similarity is treated as a result for the analysed cases only and not as evidence that RUD provisions are generally applicable to irregular buildings.

The findings can be useful to structural engineers, designers, researchers, and building authorities. If selected irregular configurations demonstrate satisfactory performance, the results can provide a basis for further investigation into extending the practical use of RUD provisions. If certain configurations show critical behaviour, the results can help identify cases where detailed structural analysis and project-specific design should be required.

The study therefore does not intend to suggest that all irregular buildings can be designed using RUD provisions. Rather, it aims to identify whether some selected irregular configurations can perform satisfactorily under the same simplified provisions used for regular buildings.


\section{Contribution of the Study}
The study contributes a controlled, like-for-like analytical comparison of one NBC 205:2024 Regular reference configuration and six selected plan-irregular configurations detailed using the same ready-to-use basis and evaluated under a common NBC 105:2025 demand framework. The contribution is not a proposal for automatic code extension. Rather, it is the combined use of force demand, complete drift profiles, CM--CR separation, nonlinear pushover capacity, performance-point response, final-step hinge evidence, and a study-defined Geometric Irregularity Index (GII) to examine where the response of selected irregular plans departs from the Regular reference.

A second contribution is the explicit separation between geometric description and seismic performance. Configurations with the same area-deficit GII are retained as distinct models, allowing the study to test whether equal GII values reproduce equal response. A third contribution is the evidence-based comparison of a representative Draft Final NBC 205:2012 building, the NBC 205:2024 Regular benchmark, and a separate four-storey height-extension case without mixing the different types of available analysis evidence. The resulting conclusions therefore identify practical applicability boundaries and evidence gaps rather than claiming universal approval of RUD detailing for irregular buildings.


\needspace{18\baselineskip}
\section{Scope and Limitations}
The study covers a representative Draft Final NBC 205:2012-designed building, the NBC 205:2024 regular model, selected plan-irregular models (L, L1, L2, T, U, and non-parallel), and a separate four-storey regular extension. The primary comparative analysis comprises seven configurations evaluated through five matched cases in each of two analysis directions. The separate four-storey model is treated outside this primary matched set and is not one of the five matched cases. The separate four-storey extension is examined through its own design-check results and is not used to infer a quantitative height effect. The comparison is made under the adopted NBC 105:2025 demand cases using storey shear, drift, descriptive CM--CR separation, pushover capacity, embedded performance points, final-step hinge evidence, and observed design-check failures.

The conclusions are limited to the analysed configurations and demand cases. The 2012 reference is assessed through concrete design checks rather than a matched pushover comparison. CM--CR separation is interpreted as an eccentricity measure and does not constitute a code-based torsional-irregularity classification. Final-step hinge patterns are distinguished from hinge states at the performance point. Staircase and landing-beam geometry contributes to the model description but is not taken as a separate member-adequacy check. The results therefore support further investigation and should not be interpreted as a formal extension of RUD provisions to all irregular buildings.
